\section{Method}\label{sec:method}

\textbf{Data:} We train and evaluate on a dataset we refer to as \emph{pop}. This is comprised of annotations from four sources — (1) \emph{Mcgill Billboard}~\cite{McgillBillboard}: songs randomly selected from the Billboard `Hot 100' Chart between 1958 and 1991, (2) \emph{Isophonics}~\cite{Isophonics}: songs from albums by The Beatles, Carole King and Zweieck, (3) \emph{RWC-Pop}~\cite{RWC}: pop songs with annotations available\footnote{\url{https://github.com/tmc323/Chord-Annotations}}, and (4) \emph{USPop}~\cite{USPop} songs chosen for popularity. Annotations and audio were kindly provided by % TODO: Remove anonymised authors 
[anonymised authors] in the ACR community. Much data integrity and preprocessing work was conducted to ensure the dataset was clean and aligned. 4 songs were not in the same tuning as the annotations, or had incorrect annotations and were removed. For further details, see the full work\footnote{\url{https://arxiv.org/abs/2512.22621}}. In total, this leaves 1,213 songs.

% TODO: Anonymise Full Work on arxiv

We split the data into 60/20/20 train/validation/test splits in a song-wise manner. We report results on the validation set for all experiments, except the final results where we train on the combined train and validation set and report on the test set, as is best practice in machine learning. While this is not the common practice in ACR, we feel that a sizeable validation set is important to obtain good estimates of the generalisation error, and that cross-fold validation may distort results in the literature~\cite{biascv}.

All audio signals are converted into a CQT using a hop length of $4096$ and a sampling rate of $44,100$Hz using \texttt{librosa}~\cite{librosa}. Different previous works have used different hop lengths. We thus test different hop lengths in Appendix~\ref{app:hop_lengths} and find that $4096$ attains among the best performances while maintaining low computational cost.

\textbf{Evaluation:} We use the standard weighted chord symbol recall (WCSR) metric for evaluation. However, to facilitate fair comparisons over differing frame lengths and beats, we define this using an integral over time, rather than discrete sum over frames as is standard. A formal definition is given in Appendix~\ref{app:weighted_chord_symbol_recall}. WCSR is then the percentage of time that the predicted chord label matches the ground truth chord label, where correctness can be measured by a variety of metrics. Such metrics include exact matching (or just `acc'), \texttt{root}, \texttt{third} and \texttt{seventh} as well as \texttt{mirex} which requires that two chords match in at least three pitches. This measure is the most forgiving, as it allows for missing upper extension, incorrect inversions or maj/min7 being confused with the relative min/maj respectively. All metrics are implemented in the \texttt{mir\_eval} library~\cite{mir_eval}. We also report the mean class-wise WCSR, which is the average of the WCSR for each chord class, also defined in Appendix~\ref{app:weighted_chord_symbol_recall}. This provides a measure of performance on the long tail of the chord distribution.

\textbf{CRNN:} All experiments are performed using a CRNN, as described by~\cite{StructuredTraining}. The model consists of a CNN feature extractor followed by a bidirectional GRU. We ablate on the model architecture by comparison to a CNN-only model, and varying the width and depth of the CNN and GRU with random search. We find that performance is not sensitive to hyperparameters within reasonable bounds and that those proposed by~\cite{StructuredTraining} are among the best choices. Thus, we stick with these choices for our baseline model. Results can be found in Appendix~\ref{app:crnn_ablation}. We use this CRNN as our baseline model for all experiments.

The is in place of a transformer model, due to both computational constraints and because the original `BTC' model is outperformed by a CRNN~\cite{BTC}. Further, ablation on increasing numbers of parameters in the CRNN suggest that model complexity is not the bottleneck. Finally, we seek to run many experiments to find sources of improvement, not to train the single best model. A lighter CRNN is more suited to this task.

\textbf{Training:} A training epoch consists of randomly sampling a patch of audio from each song in the training set. The length of this sample is kept as a hyperparameter, set to $10$ seconds for the majority of experiments. This length of patch was also tested in Appendix~\ref{app:crnn_ablation} and was found to have little effect on performance. For evaluation, the entire song is used as performance was found to be marginally better. When validating mid-way through training, songs are split into patches of the same length as the training patches to save on computation time. Samples in a batch are padded to the maximum length of a sample in the batch and padded frames are ignored for loss and metric calculation.

Experiments are run on GTX Titan Xs (12GB VRAM), RTX A6000s (48GB VRAM) and CPUs when GPUs were unavailable. All models are trained with the Adam optimiser~\cite{adam} with a learning rate of $0.001$ and cosine scheduling. Experiments testing learning rates and scheduling can be found in Appendix~\ref{app:learning_rate_scheduler}. Models are trained to minimise the cross entropy loss between the predicted chord and the true chord distribution. We use a batch size of 64 and train for 150 epochs. Validation part-way through training is conducted every 5 epochs. The model is saved whenever the validation loss improves. Each training run takes approximately 30 minutes of GPU time or 1 hour 30 minutes of CPU time, though this varies significantly across experiments.

We perform output smoothing with an HMM~\cite{BalanceRandomForestACR}, add a weighting to the loss function to improve performance on rare chord classes~\cite{ACRLargeVocab1}, separately target the root, third and seventh inspired by~\cite{StructuredTraining} and perform pitch augmentation as is standard in ACR~\cite{StructuredTraining,ACRLargeVocab1,chordformer}. Details and experiments replicating results from previous works on each of these improvements can be found in Appendices~\ref{app:output_smoothing}, \ref{app:weighted_loss}, \ref{app:structured_loss} and \ref{app:pitch_augmentation} respectively.

% \textbf{Smoothing:} Without smoothing, the CRNN predicts an average of $168$ transitions per song, compared with the $104$ in the ground truth. We thus smooth the outputs of the CRNN using a HMM on the output probabilities, bringing its predictions to $102$ transitions per song. While this does not increase overall performance, it does improve the outputs in that they are easier for a human to interpret since there are far fewer sharp chord transitions. Details of this smoothing and a comparison between HMMs and CRFs can be found in Appendix~\ref{app:output_smoothing}.

% \textbf{Weighted Loss:} A standard method of addressing a long-tailed distribution is to give greater weight in the loss to the under-represented classes. This has been used in ACR before~\cite{ACRLargeVocab1}. We use our own formulation. A description and results from varying the weighting can be found in Appendix~\ref{app:weighted_loss}. We use a conservative weighting that negligibly harms overall accuracies while improving performance on rarer chord classes.

% \textbf{Structured Decomposition:} Many previous works in ACR predict the root, third and seventh and the overall label of the chord separately~\cite{StructuredTraining,chordformer,discordharmony}. We employ a similar approach and investigate varying the coefficient of the structured decomposition in the loss affects performance. We find that performance increases noisily, and choose an optimal value for the remainder of experiments. Details of experiments can be found in Appendix~\ref{app:structured_loss}.

% \textbf{Pitch Augmentation:} We perform data augmentation with pitch shifting directly on the CQT as in~\cite{ACRLargeVocab1}. When a sample is drawn from the training set, it is shifted with probability $p=0.9$. The shift is measured in semitones in the set $\{-5,-4\ldots -1, 1, \ldots 6\}$ with equal probability of each shift. This results in 12 times as much training data. Convergence is still reached in 150 epochs. Shifting the CQT matrix is done by moving all items up or down by the number of bins corresponding to the number of semitones in the shift. An empirical comparison between shifting on the CQT and audio, as well as the probability of shift $p$ can be found in Appendix~\ref{app:pitch_augmentation}. 



